\documentclass{article}
\usepackage{amsmath,amssymb,amsthm}
\usepackage{natbib}
\usepackage{listings}
\usepackage{color}


\title{Reason Kernel Documentation with Proofs}
\author{Michał Stanisław Wójcik}
\date{\today}

\theoremstyle{plain}
\newtheorem{theorem}{Theorem}[subsection]
\newtheorem{proposition}[theorem]{Proposition}
\newtheorem{definition}[theorem]{Definition}
\newtheorem{corollary}[theorem]{Corollary}
\newtheorem{lemma}[theorem]{Lemma}
\newtheorem{fact}[theorem]{Fact}
\newtheorem{conjecture}[theorem]{Conjecture}
\newtheorem{question}[theorem]{Question}
\newtheorem{claim}[theorem]{Claim}
\newtheorem{example}[theorem]{Example}
\newtheorem{assumption}[theorem]{Assumption}
\newtheorem{remark}[theorem]{Remark}
\newtheorem{axiom}[theorem]{Axiom}
\newtheorem{note}[theorem]{Note}
\newtheorem{law}{Law}

\newcommand{\defeq}{\stackrel{\mathrm{def}}{=}}
\newcommand{\elev}{\text{elev}}
\newcommand{\free}{\text{free}}
\renewcommand{\axiom}{\text{Axiom}}
\renewcommand{\rule}{\text{Rule}}
\newcommand{\asm}{\text{Asm}}
\newcommand{\block}{\text{Block}}
\newcommand{\Sk}{\mathcal{S}}
\newcommand{\sk}{\mathfrak{s}}
\renewcommand{\i}{\mathfrak{i}}
\newcommand{\ii}{\hat{\mathfrak{i}}}
\newcommand{\N}{\mathbb{N}}
\newcommand{\F}{\mathcal{F}}
\newcommand{\C}{\mathcal{C}}
\renewcommand{\root}{\text{root}}
\newcommand{\final}{\text{final}}
\newcommand{\truth}{\tau}
\newcommand{\context}{\text{Ctx}}


\definecolor{dkgreen}{rgb}{0,0.6,0}
\definecolor{gray}{rgb}{0.5,0.5,0.5}
\definecolor{mauve}{rgb}{0.58,0,0.82}

\lstset{frame=tb,
  language=Python,
  aboveskip=3mm,
  belowskip=3mm,
  showstringspaces=false,
  columns=flexible,
  basicstyle={\small\ttfamily},
  numbers=none,
  numberstyle=\tiny\color{gray},
  keywordstyle=\color{blue},
  commentstyle=\color{dkgreen},
  stringstyle=\color{mauve},
  breaklines=true,
  breakatwhitespace=true,
  tabsize=3
}

\begin{document}

\maketitle

\section{Context}

\begin{definition}[\textbf{finite sequences}]
Let $A$ be an arbitrary set, then $A^*$ is a set of all finite sequences of $A$.
\end{definition}

We will often use notation in which we will denote sequences as $[a_1, \dots, a_n]$.
So if $a_i\in A$ then for $s=[a_1, \dots, a_n]$, we have $s\in A^*$. 

For a sequence $s=[a_1, \dots, a_n]$, $*s$ means just $a_1, \dots, a_n$, which allows us to write in a compact form many sequence operations. E.g. concatenation of sequences $s_1, s_2$ is $[*s_1, *s_2]$. If $a\in A$, then $[*s, a]$ is just a sequence with appended $a$, etc. 

To access $n-th$ element of a sequence we will use $0$-based indexing $a_i = s[i-1]$.
Also $s[-1] = a_n$.
We will denote a length of a sequence as $|s|$.
We will denote empty sequence as $[]$. 

\begin{definition}
Let $r_1, r_2\in \N^*$.
\begin{equation}
r_1 < r_2
\end{equation}
iff there exist $k \in N$ and $r, s\in\N^*$ such that $s\not=[]$, $r_1 = [*r, k]$ and $r_2 = [*r, *s]$ and $s[0] > k$.
\begin{equation}
r_1 \leq r_2
\end{equation}
iff $r_1 < r_2$ or $r_1 = r_2$.
\end{definition}

\begin{definition}
Let $r_1, r_2\in\N^*$ and $r_1 < r_2$.
\begin{equation}
\i_{r_1, r_2} \defeq \max \{ i < |r_1| : r_1[i] = r_2[i] \},
\end{equation}
\begin{equation}
\ii_{r_1, r_2} \defeq \i_{r_1, r_2} + 1.
\end{equation}
\end{definition}

\begin{corollary}
For any non empty $r\in N^*$, we have $[0] \leq r$. 
\end{corollary}

\begin{definition}
Let $r_1, r_2\in \N^*$.
\begin{equation}
r_1 \prec r_2
\end{equation}
iff there exists $s\in\N^*$ such that $s\not=[]$ and $r_2 = [*r_1, *s]$.
\begin{equation}
r_1 \preceq r_2
\end{equation}
iff $r_1 \prec r_2$ or $r_1 = r_2$.
\end{definition}

\begin{corollary}
$\preceq$ on $N^*$ is a locally finite poset.
\end{corollary}

\begin{corollary}
For any $r\in \N^*$, we have $[] \preceq r$.
\end{corollary}

\begin{lemma}
For any $r_1, r_2, r_3\in\N^*$ we have
\begin{equation}
r_1 \not< r_1,
\end{equation}
\begin{equation}
r_1 < r_2 \text{ and } r_2 < r_3 \implies r_1 < r_3.
\end{equation}
\end{lemma}
\begin{proof}
$r_1 \not< r_1$ follows directly from definition.
Assume that $\i_{r_1, r_2} \leq \i_{r_2, r_3}$.
Then 
$$
r_1[k] = r_2[k] = r_3[k]
$$
for all $k \leq \i_{r_1, r_2}$.
Note that
$$
r_1[\ii_{r_1, r_2}] < r_2[\ii_{r_1, r_2}] \leq r_3[\ii_{r_1, r_2}],
$$
thus $r_1 < r_3$.

Now, assume that $\i_{r_1, r_2} > \i_{r_2, r_3}$, then
$$
r_1[k] = r_2[k] = r_3[k]
$$
for all $k \leq \i_{r_2, r_3}$. Note that
$$
r_1[\ii_{r_2, r_3}] = r_2[\ii_{r_2, r_3}] < r_3[\ii_{r_2, r_3}],
$$
thus $r_1 < r_3$.   
\end{proof}

\begin{corollary}
$\leq$ on $N^*$ is a locally finite poset.
\end{corollary}

\begin{definition}
A $D$ is a context iff $D\subset \N^*$ and
there exists $r\in \N^*$ such that 
for all $x\in D$, we have
\begin{enumerate}
\item $r \prec x$,
\item $[*r, 0] \leq x$,
\item $\{y \in \N: [*r,0] \leq y < x \} \subset D$,
\item $\{y \in \N: r \prec y \prec x \} \subset D$.
\end{enumerate}
Such $r$ will be called a root of a context $D$.
\end{definition}

Note that $\emptyset$ is a context and for any $r\in\N^*$ r is a root of $\emptyset$.

\begin{lemma}
Let $D$ be a nonempty context with a root $r$, then $[*r,0]\in D$.
\end{lemma}
\begin{proof}
Since $D$ is nonempty, we have $x\in D$ and by definition $\{y \in \N: [*r,0] \leq y < x \} \subset D$, thus $[*r, 0]\in D$.
\end{proof}


\begin{proposition}
If $r_1$ and $r_2$ are roots of a context $D\not=\emptyset$, then $r_1 = r_2$. 
\end{proposition}
\begin{proof}
Since $[*r_1, 0], [*r_2, 0]\in D$ we have both $[*r_1, 0] \leq [*r_2, 0]$ and
$[*r_2, 0]\leq[*r_1, 0]$, thus $[*r_1, 0]=[*r_2, 0]$ and hence $r_1 = r_2$. 
\end{proof} 

We can then introduce a function $\root$ such that for any non empty $D$, $\root(D) = r$ is a root of a context $D$.

Also we will introduce for any nonempty $D$, $\overline{D} \defeq D \cup \root(D)$.

\begin{definition}
Let $D$ be a context and $x_0\in D$.
\begin{equation}
\context(D,x_0) \defeq \{ x\in D: x_0 \prec x\}.
\end{equation}
\end{definition}

\begin{theorem}
If $D$ is a context and $x_0\in D$, then $\context(D,x_0)$ is a context with a root $x_0$.
\end{theorem}

\begin{theorem}
Let $D$ be a context, $s\in\N^*$ and $k_0\in\N$. If $[*s, k_0]\in D$, then for
all $k \in \langle 0, k_0 \rangle$, we have $[*s,k]\in D$.
\end{theorem}
\begin{proof}
Take an $r\in\N^*$ which is a root of $D$. Take $[*s, k_0]\in D$.
Since $r\prec [*s, k_0]$, we have $r\preceq s$.
Since $[*r, 0]\in D$ and we have $[*r, 0] < [*s, k] < [*s, k_0]$ for $k \in (0, k_0)$, hence $[*s, k]\in D$. If $[*r, 0] < [*s,0]$ then $[*r, 0] < [*s,0] < [*s, k_0]$ and $[*s, 0]\in D$.
If $[*r, 0] \not< [*s,0]$ but $[*r, 0] < [*s,1]$, there must be $r=s$ and thus anyway $[*s,0]\in D$. 
\end{proof}

\begin{definition}
Let $D$ be a context.
\begin{equation}
\final(D) \defeq [*\root(D), k_0]\in D,
\end{equation}
such that for all $k\in\N$ and $k > k_0$, $[*\root(D), k]\not\in D$.
\end{definition}

\begin{definition}
$d:\N^*\to\N$ such that $d(\rho) \defeq |\rho| - 1$.
\end{definition}


\section{Reference System}

Let $\F$ be a set of all first order logic formulas with variable set $V$ and constants set $C$.

\begin{definition}
Let $\C$ be a set of all consts $\C \defeq \{c_0, c_1, c_2, \dots \}$.
\end{definition}

\begin{definition}
Let $\Sk$ be a set of all consts $\Sk \defeq \{\sk_r: r\in \N^*\}$.
\end{definition}

Let $\F^+$ be a set of all first order logic formulas with variable set $V$ and constants set $C \cup \C \cup \Sk$.

\begin{definition}
A pair $(D, \gamma)$ is called a reference system iff
\begin{enumerate}
\item $D\subset \N^*$ and $D$ is an notempty context,
\item $\gamma: \overline{D} \to \F^+$. 
\end{enumerate}
\end{definition}

\begin{definition}
A reference system $(D, \gamma)$ is a proof of $\psi\in\F^+$ with respect to $\mathcal{A}\subset F^+$ iff for $\eta = \root(D)$, $\gamma_\eta = \ulcorner \truth\to\psi \urcorner$
or $\gamma_\eta = \psi$
and for any $\rho\in \overline{D}$ all of the following conditions holds:
\begin{enumerate}
\item for any $c_{d}$ in $\gamma_\rho$, $d \leq d(\rho)$,
\item for any $\sk_{\rho\prime}$ in $\gamma_\rho$, $\rho\prime \leq \rho$.
\item if $\context(D, \rho) = \emptyset$, at least one of the following holds:
\begin{enumerate}
\item (AXM) $\gamma_\rho\in \mathcal{A}$,
\item (ASM) there exists $\eta \prec \rho$ 
such that $\gamma_\eta = \ulcorner \alpha_\eta \to \phi_\eta \urcorner$ and $\gamma_\rho=\alpha_\eta$.  
\item (MOD) there are $\eta_1, \eta_2 < \rho$, such that
\begin{equation}
\gamma_{\eta_1} = \ulcorner \gamma_{\eta_2} \to \gamma_\rho \urcorner,
\end{equation}
\item (CTV) there exists $\eta < \rho$ and $x\in V$ such that
\begin{equation}
\gamma_\eta = \gamma_\rho(x/c_d),
\end{equation}
\item (SKO) there exists $\eta < \rho$ and $x\in V$ and $\beta\in \F^+$ such that
$\gamma_\rho = \beta(x/sk_\rho)$ and
\begin{equation}
\gamma_\eta = \ulcorner \exists x. \beta\urcorner, 
\end{equation}
\item (GEN) there exists $\eta < \rho$ and $x\in V$ such that
\begin{equation}
\gamma_\rho = \ulcorner \forall x. \gamma_\eta \urcorner,
\end{equation}
and for any $\zeta \prec \rho$, we have $\gamma_\zeta = \ulcorner \alpha_\zeta\to\phi_\zeta\urcorner$ and $x\not\in\free(\alpha_\zeta)$,
\end{enumerate}
\item if $\context(D, \rho) \not= \emptyset$, then for 
$\eta = \final(\context(D, \rho))$, we have
\begin{equation}
\gamma_\rho = \ulcorner \alpha_\rho \to \gamma_\eta \urcorner,
\end{equation}
\end{enumerate}
\end{definition}

\section{Completness}

\begin{definition}[Axiom Schemas]
\label{axioms_schemas} 
For any formulas $\alpha, \beta, \gamma\in\F^+$ the following are logical axioms:
\begin{enumerate}
\item (THR) $\tau$
\item (LEM) $\alpha \vee \neg\alpha$
\item (IMP) $\alpha \to (\beta\to\alpha)$
\item (ANL) $\alpha \wedge \beta \to \alpha$
\item (ANR) $\alpha \wedge \beta \to \beta$
\item (AND) $\alpha \to (\beta \to (\alpha \wedge \beta))$
\item (ORL) $\alpha \to \alpha \vee \beta$
\item (ORR) $\beta \to \alpha \vee \beta$
\item (DIS) $(\alpha \to \gamma) \to ((\beta \to \gamma) \to (\alpha \vee \beta \to \gamma))$
\item (CON) $\neg \alpha \to (\alpha \to \beta)$
\item (IFI) $(\alpha \to \beta)\to((\beta\to\alpha)\to(\alpha\leftrightarrow\beta))$
\item (IFO) $(\alpha\leftrightarrow\beta) \to (\alpha\to\beta) \wedge (\beta\to\alpha)$
\item (ALL) $(\forall x. \alpha) \to \alpha(x/\tau)$ where $\tau$ is an arbitrary term and $\alpha(x/\tau)$ is admissible.
\item (EXT) $\alpha(x/\tau) \to \exists x.\alpha$ where $\tau$ is an arbitrary term and $\alpha(x/\tau)$ is admissible.
\end{enumerate}

Let $\mathfrak{A}$ denotes a set of all formulas which can be obtained by the above schemas.
\end{definition}

\begin{lemma}
Let $r\in \N^*$ and $\alpha, \beta, \gamma \in \F^+$ such that for any constant $c_d$ in them, we have $d \leq d(r)$ and for any constant $\sk_\eta$ in them, we have $\eta < r$.

There exists a reference system $(D, \gamma)$ with $\root(D) = r$ which is a proof of 
\begin{equation}
\ulcorner (\alpha \to (\beta\to\gamma))\to((\alpha\to\beta)\to(\alpha\to\gamma)) \urcorner
\end{equation}
with respect to $\mathfrak{A}$.
\end{lemma}
\begin{proof}
\begin{lstlisting}[texcl, escapechar=\%]
[*r] %$\mapsto (\alpha \to (\beta\to\gamma))\to((\alpha\to\beta)\to(\alpha\to\gamma))$%
[*r, 0] %$\mapsto (\alpha\to\beta)\to(\alpha\to\gamma)$%
[*r, 0, 0] %$\mapsto \alpha\to\gamma$%
[*r, 0, 0, 0] %$\mapsto \alpha$% # ASM $\gamma_{*r, 0, 0}$
[*r, 0, 0, 1] %$\mapsto \alpha\to\beta$% # ASM $\gamma_{*r, 0}$
[*r, 0, 0, 2] %$\mapsto \beta$% # MOD $\gamma_{*r, 0, 0, 1}, \gamma_{*r, 0, 0, 0}$
[*r, 0, 0, 3] %$\mapsto \alpha \to (\beta\to\gamma)$% # ASM $\gamma_{*r}$
[*r, 0, 0, 4] %$\mapsto \beta\to\gamma$% # MOD $\gamma_{*r, 0, 0, 3}, \gamma_{*r, 0, 0, 0, 0}$
[*r, 0, 0, 5] %$\mapsto \gamma$% # MOD $\gamma_{*r, 0, 0, 4}, \gamma_{*r, 0, 0, 0, 2}$
\end{lstlisting}
\end{proof}

\begin{lemma}[ALH]
\label{ALH}
Let $r\in \N^*$ and $\phi, \psi\in \F^+$ such that for any constant $c_d$ in them, we have $d \leq d(r)$ and for any constant $\sk_\eta$ in them, we have $\eta < r$.
If $\phi$ is a formula for which $x\not\in \text{free}(\phi)$, then there exists
a reference system $(D, \gamma)$ with $\root(D) = r$ which is a proof of 
\begin{equation}
\ulcorner \big(\forall x. (\phi\to\psi(x))\big)\to\big(\phi\to\forall x.\psi(x)\big) \urcorner
\end{equation}
with respect to $\mathfrak{A}$. 
\end{lemma}

\begin{lemma}[EXH]
\label{EXH}
Let $r\in \N^*$ and $\phi, \psi\in \F^+$ such that for any constant $c_d$ in them, we have $d \leq d(r)$ and for any constant $\sk_\eta$ in them, we have $\eta < r$.
If $\phi$ is a formula for which $x\not\in \text{free}(\phi)$, then there exists
a reference system $(D, \gamma)$ with $\root(D) = r$ which is a proof of 
\begin{equation}
\ulcorner \big(\forall x. (\psi(x)\to\phi)\big)\to\big(\exists x.\psi(x)\to\phi\big) \urcorner 
\end{equation}
with respect to $\mathfrak{A}$. 
\end{lemma}

\section{Kernel 1.0 definition}

\subsection{Preliminaries}

Notes

\begin{enumerate}
\item If not stated otherwise, we will use terminology compatible with e.g. \cite{hal2012}
\item Language is enriched by additional set of special consts $c_d$ for $d=0, 1, \dots$ and by additional set of special consts 
$sk_r$ where $r$ is an arbitrary sequence of non-negative integers.

\item $\alpha(x/\tau)$ is admissible iff no free
occurrence of $x$ in $\alpha$ is in the range of a quantifier that binds any variable which
appears in $\tau$.

\item We use a convention known from Python to denote sequences as e.g. $[0, 2, 3]$, if $r$ is a sequence of natural numbers, then by $*r$ we denote all numbers from this sequence comma separated.
So if $r=[0,1,2]$, then $[*r, 3] = [0, 1, 2, 3]$. To express that $a$ is in a sequence $r$, we will often use notation $a\in r$. 
\item For any sequence of natural numbers $r$, we define $d(r) = |r| - 1$.
\end{enumerate}

\begin{definition}
For any two sequences of non-negative integers $r_1, r_2$ we define $r_1 < r_2$ iff
there exists a sequence $r$, a non-empty sequence $s$ and an non-negative integer $k$ such that $r_1 = [*r, k]$ and $r_2 = [*r, *s]$ and $s[0] > k$.
\end{definition}

\begin{definition}
For any two sequences of nonnegative integers $r_1, r_2$ we define $r_1 \leq r_2$ iff
$r_1 < r_2$ or $r_1 = r_2$.
\end{definition}

\begin{definition}[Axiom Schemas]
\label{axioms_schemas} 
For any formulas $\alpha, \beta, \gamma$ the following are logical axioms:
\begin{enumerate}
\item (LEM) $\alpha \vee \neg\alpha$
\item (IMP) $\alpha \to (\beta\to\alpha)$
\item (ANL) $\alpha \wedge \beta \to \alpha$
\item (ANR) $\alpha \wedge \beta \to \beta$
\item (AND) $\alpha \to (\beta \to (\alpha \wedge \beta))$
\item (ORL) $\alpha \to \alpha \vee \beta$
\item (ORR) $\beta \to \alpha \vee \beta$
\item (DIS) $(\alpha \to \gamma) \to ((\beta \to \gamma) \to (\alpha \vee \beta \to \gamma))$
\item (CON) $\neg \alpha \to (\alpha \to \beta)$
\item (IFI) $(\alpha \to \beta)\to((\beta\to\alpha)\to(\alpha\leftrightarrow\beta))$
\item (IFO) $(\alpha\leftrightarrow\beta) \to (\alpha\to\beta) \wedge (\beta\to\alpha)$
\item (ALL) $(\forall x. \alpha) \to \alpha(x/\tau)$ where $\tau$ is an arbitrary term and $\alpha(x/\tau)$ is admissible.
\item (EXT) $\alpha(x/\tau) \to \exists x.\alpha$ where $\tau$ is an arbitrary term and $\alpha(x/\tau)$ is admissible.
\end{enumerate}

Let $\mathfrak{A}$ denotes a set of all formulas which can be obtained by the above schemas.
\end{definition}

\begin{definition}[Rules]
\label{rules}
We introduce the following inference rules
\begin{enumerate}
\item (MOD) $\alpha \to \beta,\; \alpha \vdash \beta$
\item (GEN) $\alpha \vdash \forall x.\alpha$
\item (CTV) $\alpha(x/c_d)\vdash \alpha$
\item (SKO) $\exists x.\alpha \vdash \alpha(x/sk_\sigma)$
\item (IDN) $\alpha \vdash \alpha$
\end{enumerate}
\end{definition}

\begin{definition}[Recursive definition of a proof step]
For any
\begin{enumerate}
\item (Reference) squence of nonnegative integers $r$,
\item (Mode) $m\in \{A, B, R, T\}$,
\item (Formula) formula $\phi$.
\end{enumerate}

we have:
\begin{enumerate}
\item An ordered quadruple $(r, m, [], \phi)$ is a \textbf{proof step}.
\item If $s$ is a sequence of proof steps, then $(r, m, s, \phi)$ is a \textbf{proof step}.
\end{enumerate}
\end{definition}

The letters $A, B, R, T$ stands for $Assumption, Block, Rule, Tautology$. It is already clear that $B$ marks the nested sequence of proof steps within a given proof step. Such a step will be called a block of proof steps or a nested block of proof steps.

\begin{definition}[Recursive definition of sub proof steps]
Let $p = (r, m, s, \phi)$ be a proof step, then 
\begin{equation}
\mathcal{S}(p) = \{p\} \cup \bigcup\{\mathcal{S}(q): q\in s\}.  
\end{equation}  
\end{definition}

\begin{definition}[Recursive definition of set of references]
Let $p = (r, m, s, \phi)$ be a proof step, then 
\begin{equation}
\mathcal{R}(p) = \{r\} \cup \bigcup\{\mathcal{R}(q): q\in s\}  
\end{equation}  
\end{definition}

When for a given proof step $p$ each sub proof steps in $\mathcal{S}(p)$ have unique references,
then we use an indexing convention under which for any refenece $\rho\in\mathcal{P}(p)$, a quadruple $(\rho, m_\rho, s_\rho, \phi_\rho)$ is a corresponding sub proof step.

\subsection{Kernel}

Reason kernel is a short program written in OCaml which checks if a block of proof steps is a proof of a formula according to the following definition. 

\begin{definition}[Recursive definition of proof]
\label{proof}
A proof step $p=(r^\prime, B, s, \phi)$ is a proof of a formula $\phi$ with respect to axioms $\mathcal{A}$ iff
$s$ = $[(r, m_r, s_r, \phi_r)]_{i=0}^k$ such that
for $i=0, \cdots, k$ we have $r = [*r^\prime, i]$ with $m_r\in \{A, B, R, T\}$ and 
\begin{enumerate}
\item for any $sk_\rho$ in formula $\phi_r$, we have $\rho \leq r$,
\item for any $c_d$ in formula $\phi_r$, we have $d \leq d(r)$,
\item for $r$ such that $m_r = A$, we must have $i=0$, $s_r=[]$ 
and $\phi_r$ is an arbitrary formula, 
such that for any $c_d$ in $\phi_r$, we have $d < d(r)$.
\item if $m_{[*r',0]} = A$, we have $\phi = \ulcorner \phi_{[*r^\prime,0]} \to \phi_{[*r^\prime, k]} \urcorner$, 
\item if $m_{[*r',0]} \not= A$, we have $\phi=\phi_{[*r^\prime, k]}$, 
\item for any $r$ such that $m_r = T$, we have $s_r=[]$ and $\phi_r\in\mathcal{A}$.
\item for any $r$ such that $m_r = R$, we have $s_r=[]$ 
and $\phi_r$ can be obtained by one of the rules from Definition \ref{rules}  
applied to chosen formulas from  
$\{f_{\rho}:\rho\in\mathcal{R}(p) \text{ and } \rho < r\}$ where
\begin{itemize}
\item for rule CTV, we have additional constrain that $d = d(r)$,
\item for rule SKO, we have additional constrain that $\sigma = r$,
\item for rule GEN, we have additional constrain that 

$x\not\in\bigcup \lbrace free(\phi_{\rho}): \rho\in\mathcal{R}(p)
\text{ and }
\rho < r
\text{ and } m_{\rho}=A \rbrace$,
\end{itemize}

\item for $m_r = B$, $(r, m_r, s_r, \phi_r)$ is a proof of $\phi_r$.
\end{enumerate}
\end{definition}

\begin{definition}
$\phi \in \text{Cn}_R(\mathcal{A})$ iff there exists a proof step $([], B, s, \phi)$ which is a proof of $\phi$ with respect to axioms $\mathcal{A}$.
\end{definition}

\subsection{Completness}

We will aim to prove that for any $\models\phi$, we have $\phi \in \text{Cn}_R(\mathfrak{A})$.

\newpage
\begin{fact}[TRN]
\label{TRN}
For formulas $\alpha, \beta, \gamma$, we have
\begin{equation}
\ulcorner (\alpha \to (\beta\to\gamma))\to((\alpha\to\beta)\to(\alpha\to\gamma)) \urcorner \in \text{Cn}_R(\mathfrak{A}).
\end{equation}
\begin{proof}
\indent
\begin{lstlisting}[texcl, escapechar=\%]
([], B, [
	([0], A, %$\alpha \to (\beta\to\gamma)$%),
	([1], B, [
		([1, 0], A, %$(\alpha\to\beta)$%),
	  	([1, 1], B, [
	  		([1, 1, 0], A, [], %$\alpha$%),
	  		([1, 1, 1], R, [], %$\beta\to\gamma$%), # MOD $\phi_{[0]}$, $\phi_{[1, 1, 0]}$
	  		([1, 1, 2], R, [], %$\beta$%), # MOD $\phi_{[1, 0]}$, $\phi_{[1, 1, 0]}$
	  		([1, 1, 3], R, [], %$\gamma$%) # MOD $\phi_{[1, 1, 1]}$, $\phi_{[1, 1, 2]}$
	  	],%$\alpha\to\gamma$%) 
	],%$((\alpha\to\beta)\to(\alpha\to\gamma)$%)
],%$(\alpha \to (\beta\to\gamma))\to((\alpha\to\beta)\to(\alpha\to\gamma)$%)
\end{lstlisting}
\end{proof}
\end{fact}

\begin{fact}[ALH]
\label{ALH}
If $\phi$ is a formula for which $x\not\in \text{free}(\phi)$ and $\psi$ is an arbitrary formula, then 
\begin{equation}
\ulcorner \big(\forall x. (\phi\to\psi(x))\big)\to\big(\phi\to\forall x.\psi(x)\big) \urcorner \in \text{Cn}_R(\mathfrak{A}).
\end{equation} 
\end{fact}
\begin{proof}
\indent
\begin{lstlisting}[texcl, escapechar=\%]
([], B, [
	([0], A, [], %$\forall x. (\phi\to\psi(x))\big)$%),
	([1], B, [], [
		([1, 0], A, [], %$\phi$%),
		([1, 1], T, [], %$\big(\forall x. (\phi\to\psi(x))\big)\to(\phi\to\psi(c_1))$%), # ALL
		([1, 2], R, [], %$\phi\to\psi(c_1)$%), # MOD $\phi_{[1, 1]}$, $\phi_{[0]}$
	  	([1, 3], R, [], %$\psi(c_1)$%), # MOD $\phi_{[1, 2]}$, $\phi_{[1,0]}$
	  	([1, 4], R, [], %$\psi(x)$%) # CTV $\phi_{[1,3]}$
	  	([1, 5], R, [], %$\forall x.\psi(x)$%) # GEN $\phi_{[1,4]}$
	],%$\phi\to\forall x.\psi(x)$%)
],%$\big(\forall x. (\phi\to\psi(x))\big)\to\big(\phi\to\forall x.\psi(x)\big)$%)
\end{lstlisting}
Note that the assumption that $x\not\in \text{free}(\phi)$ is needed for 
$\ulcorner \phi\to\psi(x) \urcorner(x/c_1)$ to become $\phi\to\psi(c_1)$ and also to apply rule GEN $\phi_{[1, 4]}$.
\end{proof}

With the above proof, it becomes more clear what is a role of a special constant $c_d$. 
It is a local universal constant of a given block of proof steps, where the index $d$ of $c_d$ is the nesting depth of the block (it is equal to the length of the reference of a given block of proof steps). Each $c_d$ can be transformed into universal quantifier according to rule CTV only in its own block due to constrain $d = d(r)$ from Definition \ref{proof}.

\begin{fact}[EXH]
\label{EXH}
If $\phi$ is a formula for which $x\not\in \text{free}(\phi)$ and $\psi$ is an arbitrary formula, then 
\begin{equation}
\ulcorner \big(\forall x. (\psi(x)\to\phi)\big)\to\big(\exists x.\psi(x)\to\phi\big) \urcorner \in \text{Cn}_R(\mathfrak{A}).
\end{equation} 
\end{fact}
\begin{proof}
\indent
\begin{lstlisting}[texcl, escapechar=\%]
([], B, [
	([0], A, [], %$\forall x. (\psi(x)\to\phi)$%),
	([1], B, [], [
		([1, 0], A, [], %$\exists x.\psi(x)$%),
		([1, 1], R [], %$\psi(sk_{[1,1]})$%), # SKO $\phi_{[1,0]}$
		([1, 2], T, [], %$\big(\forall x. (\psi(x)\to\phi)\big)\to(\psi(sk_{[1,1]})\to\phi)$%), # ALL
	  	([1, 3], R, [], %$\psi(sk_{[1,1]})\to\phi$%), # MOD $\phi_{[1, 2]}$, $\phi_{[0]}$
	  	([1, 4], R, [], %$\phi$%) # MOD $\phi_{[1,3]}$, $\phi_{[1, 1]}$.
	],%$\exists x.\psi(x)\to\phi$%)
],%$\big(\forall x. (\psi(x)\to\phi)\big)\to\big(\exists x.\psi(x)\to\phi\big)$%)
\end{lstlisting}
Note that the assumption that $x\not\in \text{free}(\phi)$ is needed for 
$\ulcorner \psi(x)\to\phi \urcorner(x/sk_{[1,1]})$ to become $\psi(sk_{[1,1})\to\phi$.
\end{proof}

\begin{definition}
Let $r_1, r_2$ be two sequences of non-negative integers.
\begin{equation}
\text{elev}(r_1, r_2) \defeq [*r_1, *r_2].
\end{equation}
\end{definition}

\begin{corollary}
Let $r, r_1, r_2$ be sequences of non-negative integers, then 
$r_1 < r_2$ iff $\text{elev}(r, r1) < \text{elev}(r, r2)$. 
\end{corollary}

\begin{definition}
Let $\phi$ be a formula and $r$ be a sequence of non-negative integers.
\begin{equation}
\text{elev}(r, \phi) = \phi\prime
\end{equation}
where $\phi\prime$ is obtain by substituting any $c_d$ from $\phi$ by $c_{d + |r|}$ and any $sk_{\rho}$ by $sk_{\text{elev}(r, \rho)}$.
\end{definition}

\begin{definition}(Recursive definition of proof step elevation)
Let $r$ be a sequence of non-negative integers
\begin{enumerate}
\item For a proof step $p=(\rho, m, [], \phi)$, 
$$\text{elev}(r, p) \defeq (\text{elev}(r, \rho), m, [], \text{elev}(r, \phi)).$$
\item For proof step $p=(\rho, m, [p_i]_{i=0}^k, \phi)$,
$$\text{elev}(r, p) \defeq (\text{elev}(r, \rho), m, [\text{elev}(r,p_i)]_{i=0}^k, \text{elev}(r, \phi)).$$
\end{enumerate}
\end{definition}

\begin{lemma}
\label{elevation-lemma}
Let $\mathcal{A}$ be a set of axioms.
Let $\eta$ be a sequence of non-negative integers.
If $\phi\in\mathcal{A}$ iff $\elev(\eta, \phi)\in\mathcal{A}$ then
a proof step $p=(r^\prime, B, s, \phi)$ is a prove of $\phi$ with respect of $\mathcal{A}$ iff
$\text{elev}(r, p)$
is a prove of $\text{elev}(r^\prime, \phi)$ with respect of $\mathcal{A}$.  
\end{lemma}
\begin{proof}
This is easy to note once you realise the following invariances:
\begin{enumerate}
\item $\rho < r$ iff $\elev(\eta, \rho) < \elev(\eta, r)$,
\item $d(\elev(\eta, r)) = d(r) + |\eta|$,
\item $\phi_r\in\mathcal{A}$ iff $\elev(\eta, \phi_r)\in\mathcal{A}$,
\item $\free(\phi_\rho) = \free(\elev(\eta, \phi_\rho))$. 
\end{enumerate}
\end{proof}

Note that $\phi_r\in\mathfrak{A}$ iff $\elev(\eta, \phi_r)\in\mathfrak{A}$

\begin{definition}
Let $\mathfrak{B}$ by all formulas which can be obtained by schemas from Fact \ref{TRN}(TRN), Fact \ref{ALH}(ALH) and Fact \ref{EXH}(EXH). Let $\mathfrak{A}^+ = \mathfrak{A} \cup \mathfrak{B}$.
\end{definition}
\begin{theorem}
If $\models \phi$, then $\phi \in \text{Cn}_R(\mathfrak{A})$
\end{theorem}
\begin{proof}
If $\models \phi$, we have $\vdash \phi$. Then by \cite{hal2012} we have a seqence of formulas $[\phi_i]_{i=1}^k$, such that either $\phi_i\in\mathfrak{A}^+$ or $\phi_i$ is obtained by rules MOD or GEN from 
$\{\phi_j: j < i\}$ for $i = 1, \dots, k$ and $\phi=\phi_k$.

Let $U = \{i: \phi_i\in \mathfrak{B}\}$. Take a sequence $u(n)$ such that
$u$ enumerates all elements from $U$ for $n=0, \dots, k^\prime$.

Let $([n], B, s_n, \phi_{u(n)})$ be a proof step which is a prove of $\phi_{u(n)}$ with respect to $\mathfrak{A}$. Such proof steps exists by Fact \ref{TRN}(TRN), Fact \ref{ALH}(ALH), Fact \ref{EXH}(EXH) and Lemma \ref{elevation-lemma}.
 
Now, consider
\begin{lstlisting}[texcl, escapechar=\%]
([], B, [
	([0], B, %$s_n$%, %$\phi_{u(0)}$%),
	...
	([%$k^\prime$%], B, %$s_{k^\prime}$%, %$\phi_{u(k^\prime)}$%),
	([%$k^\prime$% + 1], m_i, [], %$\phi_1$%),
	...
	([%$k^\prime$% + k], m_i, [], %$\phi_k$%),
], %$\phi$%)
\end{lstlisting}
where $m_i = T$ for any $\phi_i\in \mathfrak{A}$ and $m_i = R$ for any $\phi_i\not\in\mathfrak{A}$. 
In case of $i\in U$, $\phi_i$ is obtained by rule IND $\phi_{u(n)}$ where $u(n)=i$. For other $\phi_i$ we just use an original rule either MOD or GEN. Note that there is no constrain on variable $x$ in GEN since there is no assumptions in the main block of proof steps.
\end{proof}

\subsection{Soundness}
\begin{theorem}
$\text{Cn}_R(\mathfrak{A}) = \text{Cn}_R(\mathfrak{A}^+)$.
\end{theorem}
\begin{proof}
Since, we have $\mathfrak{A} \subset \mathfrak{A}^+$, then $\text{Cn}_R(\mathfrak{A}) \subset \text{Cn}_R(\mathfrak{A}^+)$. It is enough then to show $\text{Cn}_R(\mathfrak{A}^+) \subset \text{Cn}_R(\mathfrak{A})$.

Take any $\phi\in \text{Cn}_R(\mathfrak{A}^+)$. We have then a proof step $p = ([], B, s, \phi)$ which is a proof of $\phi$ with respect of $\mathfrak{A}^+$. 

Consider all $\rho\in\mathcal{R}(p)$ such that $(\rho, T, [], \phi_\rho)$ where $\phi_\rho\in\mathfrak{B}$. Then by act \ref{TRN}(TRN), Fact \ref{ALH}(ALH), Fact \ref{EXH}(EXH) and Lemma \ref{elevation-lemma}, we have 
\end{proof}

\section{Kernel 2.0 definition}

\subsection{Preliminaries}

\begin{enumerate}
\item If not stated otherwise, we will use terminology compatible with e.g. \cite{hal2012}
\item Let $\mathcal{S}$ be an infinite alphabet of skolem letters which is disjoint from consts and variable alphabets. Skolem variables will be denoted by $\mathfrak{s}_n$ for $\mathfrak{s}\in \mathcal{S}$ and $n$ non negative integers.
\item Let $\mathcal{C}$ denote infinite alphabet of consts.
\end{enumerate}

\begin{definition}(Recursive definition of a proof step). For any formula $\phi$ and consts $c\in C$ and a letter $\sk\in \Sk$
\begin{enumerate}
\item $\axiom(\phi)$  is a proof step,
\item $\rule(\phi)$ is a proof step,
\item $\asm(\phi)$ is a proof step,
\item for any sequence of proof steps $s$, $\block(c, \sk, s, \phi)$ is a proof step.  
\end{enumerate}
\end{definition}

\begin{definition}
Let $\eta=\block(c^\prime, \sk^\prime, s^\prime, \phi^\prime)$

$\prec_\eta$ is a minimal transitive relation between proof steps such that
\begin{enumerate}
\item we have $\eta \prec_\eta \sigma$ for any $\sigma\in s^\prime$, 
\item and for any step proof 
$\rho = \block(c, \sk, s, \phi)$ such that 
$\eta \prec_\eta p$, we have $p \prec_\eta \sigma$ 
for $\sigma\in s$.
\end{enumerate} 
\end{definition}

\begin{definition}
For proof steps $\rho, \sigma$ 
$$\sigma \preceq_\eta \rho$$
iff $\sigma \prec_\eta \rho$ or $\sigma = \rho$.
\end{definition}

\begin{definition}
For proof steps $\rho, \sigma$ we have 
$$\sigma <_\eta \rho$$
iff there exists $p=\block(c, \sk, s, \phi)$ such that there exist $\sigma^\prime$ such that $s = [ \dots, \sigma, \dots, \sigma^\prime, \dots]$ and either $\sigma^\prime=\rho$ or $\sigma^\prime \prec_\eta \rho$.
\end{definition}

\begin{definition}
For proof steps $\rho, \sigma$ 
$$\sigma \leq_\eta \rho$$
iff $\sigma <_\eta \rho$ or $\sigma = \rho$.
\end{definition}

\begin{definition}
Let $\eta$ be a block proof step.
$\eta \preceq_\eta \block(c, \sk, s, \phi)$ is a proof of $\phi$ with respect to a set of formulas $\mathcal{A}$ iff for any $\sigma = s[i]$ for $i = 0, \dots, |s| - 1$
\begin{enumerate}
\item if $\sigma = \asm(\phi_\sigma)$, then $i = 0$,
\item if $\sigma = \axiom(\phi_\sigma)$, then $\phi_\sigma\in\mathcal{A}$,
\item if $\sigma = \rule(\phi_\sigma)$, then either
\begin{itemize}
\item (MOD) there exist $\sigma_1, \sigma_2 <_\eta \sigma$ such that
$\phi_{\sigma_1} = \ulcorner \phi_{\sigma_2}\to \phi_\sigma \urcorner$,
\item (SKO) there exist $\sigma_1 <_\eta \sigma$ and a variable $x$ such that
$$\phi_\sigma = \alpha(x/\sk_i),$$
where $\phi_{\sigma_1} = \ulcorner \exists x.\alpha \urcorner$,
\item  
\end{itemize}
   
\end{enumerate}
\end{definition}


\bibliographystyle{alpha}
\bibliography{references}

\end{document}